大型语言模型正在进入推荐系统。它们可以读懂用户兴趣、电影标题、类型和简介，也可以帮助系统把候选项目重新排序。这样的做法让推荐结果更容易解释，也让系统能够使用自然语言资料。但是，新的能力也带来新的安全问题。传统推荐系统主要担心评分、用户画像或训练数据被改变；在使用检索增强生成的系统中，风险还会出现在检索层和索引层。如果外部文本、电影简介或元数据被放进索引，系统在检索时可能把它们当成普通资料取出。随后，排序器或大型语言模型会根据这些资料产生最终推荐。因此，攻击者不一定需要修改模型参数，也不一定需要控制后端代码，只要能让有问题的文本进入被检索的资料，就可能影响推荐列表。这个问题在大型语言模型推荐场景中更重要，因为模型会阅读候选项目的文字内容，并尝试理解其中的意思。如果文字中含有误导性的描述或类似指令的内容，模型可能把本来只是数据的文本当成需要服从的信号。本文研究的核心问题就是：当推荐系统依赖检索结果，而检索结果本身不再完全可信时，系统会怎样失效，失效又应该怎样被安全、清楚、可复现地观察。

为回答这个问题，本文设计并使用了RAGPoison Lab。它不是攻击真实平台的工具，而是一个封闭、本地、可审计的红队测试基准。系统使用一个公开电影评分数据集作为实验基础，把同一批用户、同一批电影资料和同一套评价逻辑放在受控环境中运行。为了比较正常情况和被投毒情况，系统建立两个检索索引：一个是干净索引，用来表示没有被修改的电影资料；另一个是投毒索引，用来表示经过攻击规则处理后的资料。两条路径使用相同的用户画像、查询构造、候选检索、排序流程和指标计算方式，区别只在于被检索的索引不同。这样做可以减少无关因素，让实验结果更容易解释。RAGPoison Lab还把数据准备、投毒生成、索引构建、推荐请求、排序、追踪、指标计算和报告输出分成不同步骤。这样的分层设计有两个作用。第一，它让研究者知道问题发生在哪一层：是候选项目先被检索出来，还是排序阶段把它放到更高位置，还是指标只显示了质量下降却没有显示目标暴露。第二，它让实验可以被复查。每次运行都会保存配置、指标、摘要和运行记录，方便之后检查某个结果是怎样产生的。实验文件包括运行清单、攻击配置、模型配置、指标文件、差值表和摘要文件。它们让读者可以从最后的平均分数追溯到单次运行，再从单次运行追溯到当时的检索和排序设置。前端界面和Python接口可以帮助观察结果，但本文的主要证据来自保存下来的实验文件和最终结果矩阵。这个设计也让论文能够把系统实现和实验解释连接起来，而不是只展示一组没有来源的分数。

本文比较三类攻击。第一类是定向推广，它的目标是让一个指定电影进入最终的前十推荐列表。它不一定要大幅破坏整张推荐表，只要目标项目被更多用户看见，就已经达到目的。第二类是提示注入，它把类似指令的内容放进被检索的候选文本中。本文不会展开任何可操作的攻击文本，只讨论这种风险的系统意义。提示注入最值得关注的地方，是它和大型语言模型重排相连：当候选电影的文字被送给模型，让模型按照用户兴趣重新排序时，候选文本就不只是普通说明，而是模型输入的一部分。如果模型不能稳定地区分系统指令和外部资料，注入内容就可能影响排序。第三类是非定向降质，它不追求某个目标电影，而是通过改变一部分项目内容，让整体推荐质量下降。系统也比较了两种排序思想：一种是透明的确定性排序，主要根据检索分数和类型重合度工作；另一种是大型语言模型重排，它在检索后阅读候选文本并给出顺序。这样的设计让实验可以观察不同攻击在不同流程位置上的作用，而不是只看一次模型回答是否异常。更重要的是，三类攻击对应三种不同的安全问题。定向推广更像曝光风险，用户可能在不知情的情况下看到被推高的项目；提示注入更像指令边界风险，系统需要判断哪些文本只是资料、哪些文本不能被当成命令；非定向降质更像服务质量风险，系统仍然给出推荐，但推荐列表对用户帮助变小。把这些问题分开讨论，可以避免把所有攻击都简单说成“推荐变差”。这也符合本文的红队目的：不是为了展示单一攻击技巧，而是为了让系统在可控压力下暴露不同类型的弱点。

评价部分使用HR、NDCG、MRR和ASR四个指标。HR表示前十列表中是否命中用户相关项目，NDCG关注相关项目在列表中的位置并给靠前项目更多权重，MRR关注第一个相关项目出现得早不早。三者主要衡量推荐质量。ASR表示攻击目标是否进入前十列表，主要用于有明确目标的攻击。最终实验矩阵共有十五次成功运行，没有失败记录；每类攻击有五次运行，所有结果都在前十推荐条件下计算。结果显示，不同攻击家族影响的指标并不相同。提示注入使用稠密检索和大型语言模型重排，五次运行的平均变化为：HR为\(-0.002333\)，NDCG为\(-0.001028\)，MRR为\(-0.004106\)，ASR为\(0.157370\)。这说明它对一般推荐质量的影响较小，但仍能带来明显的目标进入效果。定向推广使用混合检索和确定性排序，平均HR变化为\(-0.007210\)，NDCG为\(-0.001559\)，MRR为\(-0.005400\)，ASR为\(0.348463\)。它的主要特点是目标成功率更高，而质量下降相对有限。非定向降质同样使用混合检索和确定性排序，平均HR变化为\(-0.078049\)，NDCG为\(-0.011424\)，MRR为\(-0.033196\)，ASR不适用。它造成最大的推荐质量损失。由此可见，只看一个指标是不够的：质量指标可能看不出定向目标已经被推高，而ASR也无法解释非定向降质对整体推荐质量的伤害。十五行矩阵还说明，攻击者模型名称本身不是本文要排名的对象。更重要的是攻击家族、检索方式、排序方式和投毒强度怎样共同作用。本文因此采用谨慎解释，只把这些数值看作本地实验中的模式，而不是把某个模型说成普遍更安全或更危险。

这些结果的防御意义在于，推荐系统安全不能只检查最终答案，也不能只相信模型本身。更实际的做法是把检索资料、候选项目、排序输入和指标变化都纳入观察范围。RAGPoison Lab展示了几种有用的工程思路：保持干净索引和投毒索引分离，记录每一步的配置和输出，用多种指标同时观察目标暴露和质量下降，在把候选文本交给大型语言模型之前进行过滤或清洗，并把外部资料和系统指令清楚分开。本文并不声称已经解决投毒或提示注入问题，也不声称这些数值可以代表所有真实推荐平台。实验使用的是受控数据集、封闭环境、固定的运行矩阵和有限的模型组合，结果只能说明这个本地基准中的可观察现象。真实系统通常有更大的数据、更多业务规则、更复杂的用户行为和更多防护层，所以不能直接把本文结果扩大到所有场景。本文也没有测试所有可能的检索器、重排器、提示模板或防御策略，因此结果更适合作为风险说明和基准起点，而不是最终安全结论。即便如此，本文仍然提供了一个有价值的工程贡献：它把检索增强生成、大型语言模型推荐、数据投毒、提示注入和推荐评价放在同一个可复现实验框架中，让教师和研究者可以清楚看到风险从索引进入、经过检索和重排、最后反映到HR、NDCG、MRR和ASR上的过程。这种可审计的基准可以帮助后续研究更安全地测试防御方法，并在部署前发现推荐系统中的薄弱环节。对工程实践来说，本文最直接的提醒是：如果一个推荐系统开始让大型语言模型读取检索资料并参与排序，那么索引内容、检索规则、候选文本格式和重排提示都应被当作安全边界的一部分，而不是只在模型输出端做检查，并且这些边界需要被持续测试和记录。
